\chapter{The Language of Probability}
\label{ch:probability}

Probability theory is a precise language for reasoning about uncertainty. Its
basic objects are \emph{experiments} whose outcomes cannot be predicted with
certainty, and its goal is to assign numbers—probabilities—to the various
things that might happen. This chapter introduces that language: sample
spaces and events, the rules probabilities must obey, counting, and the
profoundly useful ideas of conditional probability, independence, and Bayes'
theorem.

\section{Sample Spaces and Events}

The set of all possible outcomes of an experiment is the \emph{sample space},
denoted $\Omega$. An \emph{event} is a subset of $\Omega$: a collection of
outcomes we might be interested in. Because events are sets, the operations of
set theory—union ($\cup$, ``or''), intersection ($\cap$, ``and''), and
complement ($A^c$, ``not'')—are exactly the operations we use to combine and
negate events.

\begin{example}
A coin is tossed three times, recording heads (H) or tails (T) each time, so
\[
  \Omega = \{HHH,\,THH,\,HTH,\,HHT,\,TTH,\,THT,\,HTT,\,TTT\}.
\]
Consider the events
$A=$``exactly two tails,''
$B=$``at least two tails,''
$C=$``no head precedes the first tail'' (equivalently, tails never appears
after a head),
and $D=$``the first toss is tails.''
Express $A,B,C,D$ as sets, and then find $A^c$, $A\cup(C\cap D)$, and
$A\cap D^c$.

\solution Reading off the outcomes,
\[
  A=\{TTH,THT,HTT\},\quad
  B=\{TTH,THT,HTT,TTT\},
\]
\[
  C=\{HHH,HHT,HTT,TTT\},\quad
  D=\{THH,TTH,THT,TTT\}.
\]
The complement of $A$ contains every outcome that is \emph{not} in $A$:
\[
  A^c=\{HHH,THH,HTH,HHT,TTT\}.
\]
Since $C\cap D=\{TTT\}$ (the only outcome in both), we have
$A\cup(C\cap D)=\{TTT,TTH,THT,HTT\}$. Finally $D^c=\{HHH,HTH,HHT,HTT\}$, and the
only outcome common to $A$ and $D^c$ is $HTT$, so $A\cap D^c=\{HTT\}$.
\end{example}

Sometimes it is the \emph{structure} of the experiment, rather than an explicit
list, that tells us what the sample space should be.

\begin{example}
A coin with $\Prob(\text{heads})=p$ is tossed repeatedly until heads appears for
the \emph{second} time. What is a natural sample space, and what is the
probability that exactly five tosses are required?

\solution The quantity of interest is the toss on which the second head lands;
this is at least $2$, so $\Omega=\{2,3,4,\dots\}$. For the second head to occur
on toss $5$, exactly one of the first four tosses is a head and the fifth is a
head. There are four positions for that early head, each sequence having
probability $p^2(1-p)^3$, so the answer is
\[
  4\,p^2(1-p)^3. \qedhere
\]
\end{example}

\section{The Rules of Probability}

A probability assignment must satisfy a few axioms: every event has probability
between $0$ and $1$, the whole sample space has probability $1$, and the
probability of a union of \emph{disjoint} events is the sum of their
probabilities. From these follow the rules we use constantly. The most
important is the \emph{addition rule} (inclusion--exclusion for two events):
\[
  \Prob(A\cup B)=\Prob(A)+\Prob(B)-\Prob(A\cap B).
\]
The subtracted term corrects for the overlap, which would otherwise be counted
twice.

\begin{example}
Suppose $\Prob(A)=\tfrac23$, $\Prob(B)=\tfrac16$, and $\Prob(A\cap B)=\tfrac19$.
Find $\Prob(A\cup B)$.

\solution By the addition rule,
\[
  \Prob(A\cup B)=\tfrac23+\tfrac16-\tfrac19
   =\tfrac{12}{18}+\tfrac{3}{18}-\tfrac{2}{18}=\tfrac{13}{18}. \qedhere
\]
\end{example}

\begin{example}
For events with $\Prob(A)=0.4$, $\Prob(B)=0.5$, and $\Prob(A\cap B)=0.1$, find
the probability that \emph{exactly one} of $A$, $B$ occurs.

\solution ``$A$ or $B$ but not both'' is the event $A\cup B$ with the overlap
removed. Since $\Prob(A\cup B)=0.4+0.5-0.1=0.8$, the probability of exactly one
is
\[
  \Prob(A\cup B)-\Prob(A\cap B)=0.8-0.1=0.7. \qedhere
\]
\end{example}

The addition rule also lets us \emph{bound} probabilities even when we do not
know the overlap exactly.

\begin{example}
Two disasters $D_1$ and $D_2$ are each rare: $\Prob(D_1)\le 10^{-6}$ and
$\Prob(D_2)\le 10^{-6}$. What can be said about the probability that at least
one occurs, and that both occur?

\solution Because probabilities are nonnegative, dropping the (subtracted)
intersection term can only increase the bound:
\[
  \Prob(D_1\cup D_2)=\Prob(D_1)+\Prob(D_2)-\Prob(D_1\cap D_2)
   \le 10^{-6}+10^{-6}=2\times10^{-6}.
\]
For the intersection, note $D_1\cap D_2\subseteq D_1$, so monotonicity gives
$\Prob(D_1\cap D_2)\le\Prob(D_1)\le 10^{-6}$.
\end{example}

\section{Counting}

When every outcome of a finite experiment is equally likely, computing a
probability reduces to counting: $\Prob(A)$ is the number of outcomes in $A$
divided by the total number of outcomes. The binomial coefficient
$\binom{n}{k}=\frac{n!}{k!\,(n-k)!}$ counts the number of ways to choose $k$
items from $n$ without regard to order.

\begin{example}
A system has five components, each either working ($1$) or failed ($0$), so an
outcome is a vector $(x_1,\dots,x_5)$. In how many ways can it happen that
exactly two components work and three have failed?

\solution We must choose which $2$ of the $5$ positions hold a $1$; the rest are
$0$. The number of such choices is $\binom{5}{2}=10$. (Explicitly, the ten
outcomes range from $(1,1,0,0,0)$ through $(0,0,0,1,1)$.)
\end{example}

\section{Conditional Probability and Independence}

The \emph{conditional probability} of $A$ given $B$ rescales probabilities to
the world in which $B$ is known to have occurred:
\[
  \Prob(A\given B)=\frac{\Prob(A\cap B)}{\Prob(B)},\qquad \Prob(B)\neq 0.
\]
Two events are \emph{independent} if knowing one tells us nothing about the
other—formally, if $\Prob(A\cap B)=\Prob(A)\Prob(B)$, equivalently
$\Prob(A\given B)=\Prob(A)$.

\begin{example}
A fair die is thrown twice. Let $A$ be ``the two throws sum to $4$'' and $B$ be
``at least one throw is a $3$.'' Compute $\Prob(A\given B)$, and decide whether
$A$ and $B$ are independent.

\solution Event $B$ consists of the eleven equally likely outcomes containing a
$3$. Of these, only $(1,3)$ and $(3,1)$ sum to $4$, so
$\Prob(A\given B)=\tfrac{2}{11}$. For independence, compare: $A=\{(1,3),(3,1),(2,2)\}$
has $\Prob(A)=\tfrac{3}{36}=\tfrac1{12}$, which is \emph{not} equal to
$\Prob(A\given B)=\tfrac2{11}$. The events are therefore dependent.
\end{example}

Independence makes ``and'' probabilities multiply, which is what lets us model
repeated trials.

\begin{example}
Two laboratory experiments are run each day in separate rooms, independently,
until at least one succeeds; each succeeds with probability $p$. What is the
probability that the first success occurs on day $n$?

\solution On any single day the probability that \emph{at least one} room
succeeds is, by the addition rule and independence,
\[
  q \eqdef \Prob(\text{success on a given day}) = p+p-p^2 = 2p-p^2.
\]
Reaching day $n$ with the first success requires $n-1$ failed days followed by a
successful one. Since days are independent,
\[
  \Prob(\text{first success on day }n) = (1-q)^{\,n-1}q
   = \bigl(1-(2p-p^2)\bigr)^{n-1}(2p-p^2). \qedhere
\]
\end{example}

\section{The Law of Total Probability and Bayes' Theorem}

Often we know the probability of an event \emph{conditional} on each of several
scenarios, and want the overall probability. The \emph{law of total
probability} stitches the scenarios together: if $B$ and $B^c$ partition the
sample space,
\[
  \Prob(A)=\Prob(A\given B)\Prob(B)+\Prob(A\given B^c)\Prob(B^c).
\]
Reversing the direction of conditioning is the job of \emph{Bayes' theorem}:
\[
  \Prob(B\given A)=\frac{\Prob(A\given B)\,\Prob(B)}{\Prob(A)}
   =\frac{\Prob(A\given B)\Prob(B)}{\Prob(A\given B)\Prob(B)+\Prob(A\given B^c)\Prob(B^c)}.
\]
This is the workhorse of medical testing, reliability, and inference. The
following examples show its range.

\begin{example}[Diagnostic test for a rare condition]
A cow is tested for BSE with sensitivity $\Prob(T\given B)=0.70$ and
false-positive rate $\Prob(T\given B^c)=0.10$, where $B$ is ``the cow is
infected'' and $T$ is ``the test is positive.'' The prevalence is
$\Prob(B)=1.3\times10^{-5}$. Find $\Prob(B\given T)$.

\solution Bayes' theorem with the law of total probability in the denominator
gives
\[
  \Prob(B\given T)=\frac{0.70\,(1.3\times10^{-5})}
   {0.70\,(1.3\times10^{-5})+0.10\,(1-1.3\times10^{-5})}\approx 9\times10^{-5}.
\]
Even after a positive test the cow is almost certainly healthy: when a
condition is extremely rare, false positives vastly outnumber true positives.
This counterintuitive phenomenon is the \emph{base-rate} effect.
\end{example}

\begin{example}[Solving for the test quality]
A breath analyzer satisfies $\Prob(A\given B)=\Prob(A^c\given B^c)=p$, where $A$
is ``the analyzer reports over the limit'' and $B$ is ``the driver is truly over
the limit.'' On a Saturday night, $5\%$ of drivers are over the limit. How
large must $p$ be so that $\Prob(B\given A)=0.9$?

\solution By Bayes' theorem,
\[
  0.9=\Prob(B\given A)=\frac{p(0.05)}{p(0.05)+(1-p)(0.95)}.
\]
Clearing the denominator,
$0.9\bigl(0.05p+0.95(1-p)\bigr)=0.05p$, which rearranges to
$\bigl(\tfrac{0.05}{0.9}+0.9\bigr)p=0.95$, giving $p\approx 0.994$. A very
accurate device is needed before a positive reading is convincing, again
because the underlying event is uncommon.
\end{example}

\begin{example}[Two independent tests]
A disease has prevalence $\Prob(D)=0.01$. A test satisfies $\Prob(T\given D)=0.98$
and $\Prob(T^c\given D^c)=0.95$. (a)~Given one positive test, find
$\Prob(D\given T)$. (b)~A second, independent test also comes back positive; now
find $\Prob(D\given T_1\cap T_2)$.

\solution (a)~Bayes' theorem yields
\[
  \Prob(D\given T)=\frac{0.98(0.01)}{0.98(0.01)+0.05(0.99)}
   =\frac{0.0098}{0.0593}\approx 0.166.
\]
(b)~The two tests are independent \emph{conditional on disease status}, so
\begin{align*}
  \Prob(D\given T_1\cap T_2)
   &=\frac{\Prob(T_1\given D)\Prob(T_2\given D)\Prob(D)}
   {\Prob(T_1\given D)\Prob(T_2\given D)\Prob(D)+\Prob(T_1\given D^c)\Prob(T_2\given D^c)\Prob(D^c)}\\[4pt]
   &=\frac{(0.98)^2(0.01)}{(0.98)^2(0.01)+(0.05)^2(0.99)}\approx 0.795.
\end{align*}
\emph{A subtle point.} It is tempting to write
$\Prob(T_1\cap T_2)=\Prob(T_1)\Prob(T_2)$, but that is wrong: the tests are
\emph{not} unconditionally independent. They are independent only once disease
status is fixed. Indeed a first positive raises the chance of disease, which in
turn raises the chance the second test is positive. This distinction—between
independence and \emph{conditional} independence—is one of the most important
in all of probability.
\end{example}

\begin{example}[A multiple-choice exam]
On each question a student either knows the answer (then answers correctly with
probability $1$) or guesses among the options (then correct with probability
$\tfrac14$). The student knows each answer with probability $0.6$. Given that a
question is answered correctly, what is the probability the student actually
knew it?

\solution Let $B=$``knows the answer,'' $A=$``answers correctly.'' Then
$\Prob(B)=0.6$, $\Prob(A\given B)=1$, $\Prob(A\given B^c)=\tfrac14$, and Bayes'
theorem gives
\[
  \Prob(B\given A)=\frac{1\cdot 0.6}{1\cdot 0.6 + \tfrac14\cdot 0.4}
   =\frac{0.6}{0.7}=\frac{6}{7}. \qedhere
\]
\end{example}

\begin{example}[The neighbor and the plant]
While you are away, a neighbor waters a sickly plant. Without water it dies with
probability $0.8$; with water it dies with probability $0.15$. You are $90\%$
sure the neighbor remembers. (a)~Find the probability the plant is alive on your
return. (b)~If it is dead, find the probability the neighbor forgot.

\solution Let $A=$``alive'' and $R=$``remembers to water.'' By the law of total
probability,
\[
  \Prob(A)=\Prob(A\given R)\Prob(R)+\Prob(A\given R^c)\Prob(R^c)
   =(0.85)(0.9)+(0.2)(0.1)=\tfrac{157}{200}.
\]
For part (b), Bayes' theorem with $A^c$ gives
\[
  \Prob(R^c\given A^c)=\frac{\Prob(A^c\given R^c)\Prob(R^c)}{\Prob(A^c)}
   =\frac{(0.8)(0.1)}{1-\tfrac{157}{200}}=\frac{0.08}{43/200}=\frac{16}{43}. \qedhere
\]
\end{example}

\begin{example}[Inferring an infestation from a union]
Grapefruit is grown in two regions, $R_1$ and $R_2$. Let $A$ and $B$ be the
events that $R_1$ and $R_2$ respectively are infested with a parasite, with
$\Prob(A)=\tfrac34$, $\Prob(B)=\tfrac25$, and $\Prob(A\cup B)=\tfrac45$. If a
shipment from $R_1$ is found infested, what is the probability $R_2$ is infested
too?

\solution We want $\Prob(B\given A)=\Prob(A\cap B)/\Prob(A)$. The addition rule
recovers the intersection:
$\Prob(A\cap B)=\Prob(A)+\Prob(B)-\Prob(A\cup B)=\tfrac34+\tfrac25-\tfrac45=\tfrac{7}{20}$.
Hence $\Prob(B\given A)=\tfrac{7/20}{3/4}=\tfrac{7}{15}$.
\end{example}

\subsection{Bayes' theorem with geometric structure: the three-card problem}

\begin{example}[Three cards]
Three cards are placed in a hat: one green on both sides, one red on both sides,
and one green on one side and red on the other. You draw a card and look at one
side, both chosen at random. The side you see is green. What is the probability
the other side is also green?

\solution Let $G_1$ be ``the visible side is green'' and $G_2$ be ``the hidden
side is green.'' The only way \emph{both} sides are green is the green--green
card, so $\Prob(G_1\cap G_2)=\tfrac13$. By symmetry $\Prob(G_1)=\tfrac12$ (half
of all six card-faces are green). Therefore
\[
  \Prob(G_2\given G_1)=\frac{\Prob(G_1\cap G_2)}{\Prob(G_1)}
   =\frac{1/3}{1/2}=\frac{2}{3}.
\]
The intuitive guess of $\tfrac12$ is wrong: seeing a green face makes the
double-green card more likely than the mixed card, because the double-green card
has \emph{two} ways to show you a green face.
\end{example}

\subsection{The Monty Hall problem}

\begin{example}[Monty Hall]
A car sits behind one of three doors; goats sit behind the other two. You pick a
door, and the host—who knows where the car is—opens a different door to reveal a
goat. Should you switch? Label the door hiding the car $a$ and the others $b,c$,
and model the experiment by the product space $\Omega=\{a,b,c\}\times\{a,b,c\}$,
where the first coordinate is your choice and the second is the door the host
opens.

\solution The host never opens door $a$ (the car is there), so any pair with
second coordinate $a$ has probability $0$. If you pick $a$ (probability
$\tfrac13$), the host chooses $b$ or $c$ equally, giving
$\Prob((a,b))=\Prob((a,c))=\tfrac13\cdot\tfrac12=\tfrac16$. If you pick $b$, the
host must open $c$, so $\Prob((b,c))=\tfrac13$; similarly $\Prob((c,b))=\tfrac13$.

A \emph{no-switch} player wins exactly when the original pick is $a$, an event of
probability $\tfrac16+\tfrac16=\tfrac13$. A \emph{switching} player wins exactly
when the original pick was $b$ or $c$, namely on $\{(b,c),(c,b)\}$, of
probability $\tfrac13+\tfrac13=\tfrac23$. Switching doubles your chances.
\end{example}

\begin{center}
\begin{tikzpicture}[font=\small,>=Stealth, level distance=15mm,
   level 1/.style={sibling distance=34mm},
   level 2/.style={sibling distance=15mm}]
\node[circle,draw=accent,fill=accentlight,inner sep=1.5pt] {pick}
  child {node[draw,rounded corners,fill=softgrey] {pick $a$ ($\tfrac13$)}
    child {node {host $b$} edge from parent node[left,font=\tiny]{$\tfrac12$}}
    child {node {host $c$} edge from parent node[right,font=\tiny]{$\tfrac12$}}}
  child {node[draw,rounded corners,fill=softgrey] {pick $b$ ($\tfrac13$)}
    child {node {host $c$} edge from parent node[right,font=\tiny]{$1$}}}
  child {node[draw,rounded corners,fill=softgrey] {pick $c$ ($\tfrac13$)}
    child {node {host $b$} edge from parent node[right,font=\tiny]{$1$}}};
\end{tikzpicture}
\captionof{figure}{The Monty Hall experiment. Staying wins only on the left
branch (probability $\tfrac13$); switching wins on the other two
(probability $\tfrac23$).}
\end{center}
